\documentclass[11pt,a4paper]{article}
\usepackage[utf8]{inputenc}
\usepackage[margin=1in]{geometry}
\usepackage{graphicx}
\usepackage{amsmath,amssymb,amsfonts}
\usepackage{booktabs}
\usepackage{float}
\usepackage{hyperref}
\usepackage{microtype}
\usepackage{fancyhdr}
\usepackage{xcolor}

\pagestyle{fancy}
\fancyhf{}
\rhead{Phase II: Surgical RCM Simulation}
\lhead{Robotic Systems in Manipulation}
\rfoot{Page \thepage}

\hypersetup{
    colorlinks=true,
    linkcolor=blue,
    filecolor=magenta,      
    urlcolor=cyan,
    pdftitle={Phase II: Surgical RCM Control and Trajectory Verification},
    pdfpagemode=FullScreen,
}

\title{Surgical RCM Control and Trajectory Verification\\Phase II Project Delivery}
\author{Rodrigo Sales Freire dos Santos (ist1119658) \\ Mikhail Mohamed Shueb (ist1104014)}
\date{June 2026}

\begin{document}

\maketitle

\begin{abstract}
This report details the implementation, tuning, and verification of the Remote Center of Motion (RCM) controller for a 7-DOF KUKA LBR MED robot carrying a $15\text{ cm}$ biopsy needle tool. Designed specifically for minimally invasive surgery, the control system enforces a strict spatial constraint at the patient's entry incision (the trocar), while simultaneously tracking a multi-point 3D target path inside the body. By tuning the Closed-Loop Inverse Kinematics (CLIK) controller gains and damping factors, we achieved sub-millimeter tracking accuracy ($<0.19\text{ mm}$ peak trajectory error). The mathematical derivations, controller architecture, and simulation results are explained in detailed, beginner-friendly terms.
\end{abstract}

\section{Introduction to Remote Center of Motion (RCM)}
In minimally invasive surgery (MIS), the robot inserts a long, rigid surgical tool (like a needle or laparoscope) into the patient's body through a tiny skin incision called a \textbf{trocar}. 

To prevent tearing the patient's skin and causing tissue damage, the tool shaft must remain perfectly centered in this incision point. This physical constraint is called a **Remote Center of Motion (RCM)**. It is ``remote'' because the center of rotation is not located at any physical joint of the robot arm; instead, it is a virtual point in space suspended along the needle shaft.

In control terms, this means the robot is restricted: it cannot translate the needle shaft sideways at the trocar. It can only slide the needle forward/backward (insertion depth) and pitch/yaw the needle around the trocar center. This report documents how we control the 7-DOF KUKA LBR MED robot to accomplish this task.

\section{Workspace Setup and Target Trajectory}

\subsection{Needle and Trocar Parameters}
The physical parameters of our simulation are defined as:
\begin{itemize}
    \item \textbf{Needle Length ($L$):} $0.15\text{ m}$ ($15\text{ cm}$).
    \item \textbf{Trocar Position ($c_r$):} A fixed point in the robot's workspace coordinates:
    \begin{equation}
        c_r = \begin{bmatrix} -0.5 \\ 0.0 \\ 0.4 \end{bmatrix} \text{ meters}
    \end{equation}
\end{itemize}

\subsection{Optimized 3D Target Path}
The trajectory is planned across 6 spatial targets ($m_0$ to $m_5$) inside the patient. These points were carefully selected to cover a rich 3D workspace (varying in depth and lateral width) while remaining within the reachable limits of the $15\text{ cm}$ needle:
\begin{itemize}
    \item $m_0 = [-0.50; \phantom{-}0.00; 0.30]\text{ m}$ (Starting point)
    \item $m_1 = [-0.50; \phantom{-}0.08; 0.28]\text{ m}$
    \item $m_2 = [-0.52; \phantom{-}0.04; 0.27]\text{ m}$ (Deepest insertion target)
    \item $m_3 = [-0.46; -0.04; 0.32]\text{ m}$ (Highest/shallowest target)
    \item $m_4 = [-0.50; -0.08; 0.28]\text{ m}$
    \item $m_5 = [-0.55; \phantom{-}0.00; 0.30]\text{ m}$ (Final target and hold)
\end{itemize}

The simulation runs for 30 seconds total, transitioning smoothly from target to target every 5 seconds, followed by a static 5-second hold at $m_5$.

\section{Kinematics and Constraints (Simple Explanation)}

\subsection{Calculating the Needle Tip Position}
The robot's hand (flange) is located at position $p_e$, and its orientation is represented by a rotation matrix $R_e$. The third column of $R_e$ represents the local Z-axis direction of the hand, which we define as the direction vector $\boldsymbol{z}_e$.

Since the needle is straight and attached to the hand, its tip position ($p_{tool}$) is simply:
\begin{equation}
    p_{tool} = p_e + L \boldsymbol{z}_e
\end{equation}
This means: ``start at the hand position, and walk $15\text{ cm}$ in the direction the hand is pointing.''

\subsection{Formulating the Trocar Constraint Error}
At any instant, the trocar center $c_r$ must lie on the needle shaft. We calculate the distance the needle is inserted through the trocar (the insertion depth $\lambda$) by projecting the trocar vector onto the needle direction vector:
\begin{equation}
    \lambda = (c_r - p_e)^T \boldsymbol{z}_e
\end{equation}
The closest point on the needle shaft to the trocar is:
\begin{equation}
    p_c = p_e + \lambda \boldsymbol{z}_e
\end{equation}
The distance between the trocar center and our needle shaft is the RCM constraint error:
\begin{equation}
    e_{rcm} = c_r - p_c
\end{equation}
Our primary control goal is to keep this error magnitude $\|e_{rcm}\|$ at zero.

\section{Closed-Loop Inverse Kinematics (CLIK) Controller}

\subsection{Task-Priority Framework}
Since the robot has 7 joints (redundancy) but our task space has fewer dimensions, we can prioritize our objectives:
\begin{enumerate}
    \item \textbf{Primary Task (Safety):} Keep the trocar constraint error at zero.
    \item \textbf{Secondary Task (Tracking):} Move the needle tip to follow the trajectory.
\end{enumerate}

To do this, we calculate the joint velocities $\dot{q}$ in two parts:
\begin{equation}
    \dot{q} = \dot{q}_{rcm} + \dot{q}_{tooltip}
\end{equation}

The primary velocity $\dot{q}_{rcm}$ drives the trocar error to zero:
\begin{equation}
    \dot{q}_{rcm} = J_{rcm}^{\dagger} (K_{rcm} e_{rcm})
\end{equation}
where $J_{rcm}^{\dagger}$ is the pseudoinverse of the constraint Jacobian and $K_{rcm}$ is the safety gain.

The secondary velocity $\dot{q}_{tooltip}$ tracks the target path. Crucially, we filter this velocity through a **Null-Space Projection Matrix** $(I - J_{rcm}^{\dagger} J_{rcm})$:
\begin{equation}
    \dot{q}_{tooltip} = (I - J_{rcm}^{\dagger} J_{rcm}) J_{tool\_projected}^{\dagger} \boldsymbol{v}_{required}
\end{equation}
The Null-Space Projection acts as a ``mathematical shield'': it filters out any joint motion that would disturb the trocar. This guarantees that the secondary tracking job never violates the primary safety job.

\subsection{Controller Parameter Tuning}
During initial testing with standard parameters ($K_{tool} = 5$, $\lambda_{damp} = 0.1$), the robot suffered from significant tracking lag (up to $13.5\text{ mm}$ error) because the controller was overly cautious. We resolved this by tuning two parameters:
\begin{itemize}
    \item \textbf{Increased Tooltip Gain ($K_{tool} = 20$):} This acts like a tighter spring pulling the needle tip toward the target, forcing the controller to correct errors much faster.
    \item \textbf{Decreased Damping ($\lambda_{damp} = 0.01$):} Damping acts like a brake pedal to keep joint velocities stable near singular configurations. By lowering it from $0.1$ to $0.01$, we released the ``brakes,'' allowing the robot to move freely and follow the fast dynamic path.
\end{itemize}

\section{Simulation Results and Verification}
The simulation was executed for 30 seconds. The tracking results are highly precise and fully sub-millimeter.

\subsection{Visual Path Analysis}
The 3D trajectory plot (Fig.~\ref{fig:trajectory_3d}) demonstrates that the needle shaft (orange lines) consistently enters through the trocar (cyan triangle) and tracks the targets (green spheres). To allow visual verification, the planned path (red dashed line) is shifted by $+3\text{ mm}$ along the Y-axis so it runs side-by-side with the blue executed path; they remain parallel throughout the motion, indicating negligible tracking error.

\subsection{Trocar Constraint Precision}
Figure~\ref{fig:rcm_error} shows the trocar constraint error magnitude over time. The maximum transient trocar error is capped at $0.105\text{ mm}$ and settles to $0.000\text{ mm}$ at steady-state, ensuring absolute patient safety.

\subsection{Error Margin Statistics}
The calculated errors at target points are shown in Table~\ref{tab:target_errors}, and transition errors are shown in Table~\ref{tab:transition_errors}. 

\begin{table}[H]
\centering
\caption{Error margin at each target point.}
\label{tab:target_errors}
\begin{tabular}{@{}cccc@{}}
\toprule
Target Point & Time (s) & Trocar RCM Error (mm) & Tooltip Tracking Error (mm) \\ \midrule
$m_0$ (Start) & 0  & 0.007143 & 0.011697 \\
$m_1$         & 5  & 0.008550 & 0.065520 \\
$m_2$         & 10 & 0.000953 & 0.129795 \\
$m_3$         & 15 & 0.003167 & 0.049255 \\
$m_4$         & 20 & 0.022055 & 0.141048 \\
$m_5$ (Start) & 25 & 0.006190 & 0.036866 \\
$m_5$ (End)   & 30 & 0.000000 & 0.000000 \\ \bottomrule
\end{tabular}
\end{table}

\begin{table}[H]
\centering
\caption{Average and maximum tracking errors during transitions.}
\label{tab:transition_errors}
\begin{tabular}{@{}cccccc@{}}
\toprule
Transition & Interval (s) & Avg RCM Err (mm) & Max RCM Err (mm) & Avg Tooltip Err (mm) & Max Tooltip Err (mm) \\ \midrule
$m_0 \rightarrow m_1$ & $[0, 5]$   & 0.009865 & 0.105213 & 0.020335 & 0.185847 \\
$m_1 \rightarrow m_2$ & $[5, 10]$  & 0.003946 & 0.010926 & 0.028859 & 0.129795 \\
$m_2 \rightarrow m_3$ & $[10, 15]$ & 0.022614 & 0.098094 & 0.038649 & 0.129795 \\
$m_3 \rightarrow m_4$ & $[15, 20]$ & 0.012369 & 0.049387 & 0.034262 & 0.141048 \\
$m_4 \rightarrow m_5$ & $[20, 25]$ & 0.013092 & 0.035113 & 0.034027 & 0.141048 \\
$m_5$ Hold            & $[25, 30]$ & 0.001101 & 0.006190 & 0.007714 & 0.036866 \\ \bottomrule
\end{tabular}
\end{table}

\begin{figure}[H]
\centering
\includegraphics[width=0.75\textwidth]{images/step6_rcm_trajectory.png}
\caption{3D simulation of the tooltip trajectory and needle shaft passing through the trocar from above.}
\label{fig:trajectory_3d}
\end{figure}

\begin{figure}[H]
\centering
\includegraphics[width=0.75\textwidth]{images/step6_rcm_error.png}
\caption{Trocar constraint error magnitude over time.}
\label{fig:rcm_error}
\end{figure}

\subsection{Note on initial $t = 0$ micro-error}
A microscopic error is visible at the very start ($0.011\text{ mm}$ for tooltip and $0.007\text{ mm}$ for trocar). This is an artifact of the Simulink block parameters: the initial joint positions are entered as a rounded text string (`[0; 0.4085; 3.1416; ...]`) which introduces a tiny truncation discrepancy. This is physically negligible (equivalent to $1/10$th of a human hair width) and is corrected immediately as the simulation starts.

\section{Conclusion}
Phase II delivery successfully demonstrates a safe, highly accurate RCM controller for minimally invasive surgical procedures. By tuning the CLIK gain ($K_{tool} = 20$) and dampening ($\lambda_{damp} = 0.01$), the tracking errors were successfully reduced to a sub-millimeter level ($< 0.19\text{ mm}$) while maintaining trocar constraint safety.

\end{document}
